% !TeX root = ../main.tex
% Fonte: Lezione 7, pp. 1--3. Generalizzazioni esplicitate nel testo.

\section{Somme di sottospazi e somme dirette}

\subsection{Il sottospazio somma}

Per ottenere un sottospazio che contenga $U$ e $W$, si considera il sottospazio generato dalla loro unione, che può non essere chiusa per somma.

\begin{definizione}
La \emph{somma} di due sottospazi $U,W$ di uno spazio vettoriale $V$ è
\[
  U+W=\langle U\cup W\rangle.
\]
\end{definizione}

\begin{teorema}
Per due sottospazi $U,W\subseteq V$ si ha
\[
  U+W=\{u+w:u\in U,\ w\in W\}.
\]
\end{teorema}
\begin{proof}
Indichiamo con $T$ l'insieme delle somme a destra.

\passomath{U+W\subseteq T} Un elemento di $\langle U\cup W\rangle$ è una combinazione lineare finita di vettori di $U\cup W$. Raggruppando i termini secondo il sottospazio di provenienza, si scrive
\[
  \lambda_1u_1+\cdots+\lambda_ru_r
  +\mu_1w_1+\cdots+\mu_sw_s.
\]
La prima somma appartiene a $U$ e la seconda a $W$, per la chiusura rispetto alle combinazioni lineari. Un gruppo privo di termini dà $0_V$. Il vettore ha quindi la forma $u+w$ e appartiene a $T$.

\passomath{T\subseteq U+W} Se $u\in U$ e $w\in W$, allora $u+w=1_Ku+1_Kw$ è una combinazione lineare di elementi di $U\cup W$ e appartiene a $\langle U\cup W\rangle=U+W$.

Le due inclusioni concludono la prova.
\end{proof}

\begin{osservazione}[Verifica diretta e minimalità]
L'insieme delle somme contiene $0_V=0_V+0_V$ ed è chiuso per le operazioni:
\begin{align*}
  (u_1+w_1)+(u_2+w_2)&=(u_1+u_2)+(w_1+w_2)\in U+W,\\
  \lambda(u_1+w_1)&=\lambda u_1+\lambda w_1\in U+W.
\end{align*}
Inoltre $u=u+0_V$ e $w=0_V+w$, dunque contiene entrambi i sottospazi. Ogni sottospazio che contiene $U$ e $W$ contiene tutte queste somme: $U+W$ è il più piccolo sottospazio contenente $U\cup W$.
\end{osservazione}

\begin{esempio}[Somma dei due assi]
In $\R^2$, $(a,b)=(a,0)+(0,b)$: la somma degli assi coordinati è tutto il piano, mentre l'unione comprende soltanto i punti dei due assi.
\end{esempio}

\textbf{Più sottospazi.} Per un numero finito di sottospazi,
\[
  U_1+\cdots+U_r
  =\langle U_1\cup\cdots\cup U_r\rangle
  =\{u_1+\cdots+u_r:u_i\in U_i\}.
\]
La prova usa lo stesso raggruppamento dei termini. L'ordine dei sottospazi e la disposizione delle parentesi non cambiano il risultato, per commutatività e associatività della somma in $V$.

\subsection{Unicità della scomposizione e somma diretta}

L'appartenenza a $U+W$ assicura una scomposizione, che può non essere unica. Se $U=W$ è una retta non nulla e $v\ne0_V$ è un suo vettore, le scritture $v=v+0_V$ e $v=0_V+v$ hanno componenti diverse nelle due posizioni.

\begin{definizione}
La somma $U_1+\cdots+U_r$ si dice \emph{diretta} se ogni suo vettore si scrive in modo unico come $u_1+\cdots+u_r$, con $u_i\in U_i$. Si usa allora la notazione $U_1\oplus\cdots\oplus U_r$.
\end{definizione}

\begin{teorema}
La somma di due sottospazi $U$ e $W$ è diretta se e solo se $U\cap W=\{0_V\}$.
\end{teorema}
\begin{proof}
\passomath{\Rightarrow} Supponiamo che la somma sia diretta. Per $z\in U\cap W$, le scritture $z=z+0_V$ e $z=0_V+z$ sono entrambe scomposizioni valide. L'unicità delle componenti in $U$ impone $z=0_V$. Poiché il vettore nullo appartiene a entrambi i sottospazi, $U\cap W=\{0_V\}$.

\passomath{\Leftarrow} Supponiamo $U\cap W=\{0_V\}$. L'esistenza della scomposizione segue dalla caratterizzazione di $U+W$. Se $u+w=u'+w'$, con $u,u'\in U$ e $w,w'\in W$, allora
\[
  u-u'=w'-w\in U\cap W.
\]
Entrambi i membri sono quindi nulli: $u=u'$ e $w=w'$. La scomposizione è unica e la somma è diretta.
\end{proof}

\begin{esercizio}[Dimostrazione alternativa per assurdo]
Dimostrare che, se $U\cap W$ contiene un vettore non nullo, la scomposizione dei vettori di $U+W$ non è unica.
\end{esercizio}
\begin{proof}[Svolgimento]
Scegliamo $z\in U\cap W$, con $z\ne0_V$. A ogni scomposizione $v=u+w$ si affianca
\[
  v=(u+z)+(w-z).
\]
Le nuove componenti appartengono ancora rispettivamente a $U$ e a $W$, ma sono diverse dalle precedenti: $u+z\ne u$ e $w-z\ne w$. Ogni vettore della somma ha quindi due scomposizioni distinte. Assumere contemporaneamente la somma diretta e un'intersezione non nulla porta a una contraddizione.
\end{proof}

\Needspace{12\baselineskip}
\begin{osservazione}[Criterio per più sottospazi]
La somma $U_1+\cdots+U_r$ è diretta se e solo se
\[
  u_1+\cdots+u_r=0_V,\quad u_i\in U_i
  \quad\Longrightarrow\quad u_1=\cdots=u_r=0_V.
\]
\end{osservazione}
\begin{proof}
\passomath{\Rightarrow} Si confronta la scomposizione nulla assegnata con $0_V=0_V+\cdots+0_V$. L'unicità impone che tutte le componenti siano nulle.

\passomath{\Leftarrow} Sottraendo due scomposizioni dello stesso vettore si ottiene una somma nulla, con ciascuna componente nel rispettivo sottospazio. La proprietà impone che tutte le differenze siano nulle, quindi le scomposizioni coincidono.
\end{proof}

\textbf{Attenzione alle intersezioni.} Per tre o più sottospazi, le sole intersezioni a due a due nulle non garantiscono la somma diretta.\footnote{In $\R^2$, le rette generate da $(1,0)$, $(0,1)$ e $(1,1)$ hanno intersezioni a due a due nulle, ma $(1,0)+(0,1)-(1,1)=0$ è una relazione con componenti non nulle. Per due sottospazi, invece, l'intersezione nulla è il criterio necessario e sufficiente.}
